%# -*- coding: utf-8-unix -*-
%==================================================
%% bjtuthesis.cls for BJTU Master Thesis
%%==================================================

\NeedsTeXFormat{LaTeX2e}[1995/12/01]
\ProvidesClass{bjtuthesis}[2015/3/15 v0.9 Beijing Jiaotong University Template]
\def\version{v0.9}
\newif\ifbjtu@bachelor\bjtu@bachelorfalse
\newif\ifbjtu@master\bjtu@masterfalse
\newif\ifbjtu@doctor\bjtu@doctorfalse
\newif\ifbjtu@review\bjtu@reviewfalse
\newif\ifbjtu@submit\bjtu@submitfalse
\newif\ifbjtu@professional\bjtu@professionalfalse
\DeclareOption{bachelor}{\bjtu@bachelortrue}
\DeclareOption{master}{\bjtu@mastertrue}
\DeclareOption{doctor}{\bjtu@doctortrue}
\DeclareOption{review}{\bjtu@reviewtrue}
\DeclareOption{submit}{\bjtu@submittrue}
\DeclareOption{professional}{\bjtu@professionaltrue}
\DeclareOption*{\PassOptionsToClass{\CurrentOption}{ctexbook}}
\ProcessOptions\relax
\LoadClass[a4paper,UTF8,scheme=chinese]{ctexbook}
\ifbjtu@bachelor\relax\else
  \ifbjtu@master\relax\else
    \ifbjtu@doctor\relax\else
      \ClassError{bjtuthesis}%
      		{You have to specify one of thesis options: bachelor, master or doctor.}{}
    \fi
  \fi
\fi

%% bjtuthesis.cls segments
% 0. Import bjtuthesis.cfg
% 1. Import and configure LaTeX packages.
% 2. Define general-purpose LaTeX commands. 
% 3. Configure the imported packages, also extend LaTeX command in bjtuthesis
% 4. Draw the bjtuthesis

%==========
% Segment 0. Import bjtuthesis.cfg 
%==========

%% 导入 bjtuthesis.cfg 文件
\AtEndOfClass{%# -*- coding: utf-8-unix -*-
%%==================================================
%% bjtuthesis.cfg for BJTU Thesis
%%==================================================

\ProvidesFile{bjtuthesis.cfg}[2015/05/01 v0.8 bjtuthesis configuration file]

%% 目录、插图索引、表格索引
\def\bjtu@contentsname{目录}
\def\bjtu@listfigurename{插图索引}
\def\bjtu@listtablename{表格索引}
\def\bjtu@listalgorithmname{算法索引}

%%
%% labels in the title page
%%
\def\bjtu@label@major{专业}
\def\bjtu@label@title{论文题目}
\def\bjtu@label@author{论文作者}
\def\bjtu@label@thesis{学位论文}
\def\bjtu@label@advisor{导师}
\def\bjtu@label@advisorjobtitle{职称}
\def\bjtu@label@coadvisor{副导师}
\def\bjtu@label@defenddate{答辩日期}
\def\bjtu@label@studentnumber{学号}

%%
%% string values filled in the title page
%%

\def\bjtu@value@classification{}
\def\bjtu@value@confidential{}
\def\bjtu@value@serialnumber{}
\def\bjtu@value@school{}
\ifbjtu@bachelor
  \def\bjtu@value@chinesedegree{学士}
  \def\bjtu@value@englishdegree{Bachelor}
\else
  \ifbjtu@master
    \def\bjtu@value@chinesedegree{硕士}
    \def\bjtu@value@englishdegree{Master}
  \else
    \ifbjtu@doctor
      \def\bjtu@value@chinesedegree{博士}
    \def\bjtu@value@englishdegree{Doctor}
    \else
      \ClassError{bjtuthesis}%
      {Unknown value for degree.}{}
    \fi 
  \fi
\fi

\ifbjtu@professional\relax
  \def\bjtu@label@statement{\bjtu@value@chinesedegree{专业}\bjtu@label@thesis}
\else
  \def\bjtu@label@statement{\bjtu@value@chinesedegree\bjtu@label@thesis}
\fi
% 论文独创性声明
\def\bjtu@label@original{独创性声明}
\def\bjtu@label@authorization{学位论文版权使用授权书}
\def\bjtu@label@authorsign{学位论文作者签名：}
\def\bjtu@label@Supervisorsign{导师签名：}
\def\bjtu@label@originalDate{签字日期：\hfill\hfill 年 \hfill 月 \hfill 日}
\def\bjtu@label@originalcontent{本人声明所呈交的学位论文是本人在导师指导下进行的研究工作和取得的研究成果，除了文中特别加以标注和致谢之处外，论文中不包含其他人已经发表或撰写过的研究成果，也不包含为获得北京交通大学或其他教育机构的学位或证书而使用过的材料。与我一同工作的同志对本研究所做的任何贡献均已在论文中作了明确的说明并表示了谢意。
}
\def\bjtu@label@authorizationcontent{本学位论文作者完全了解北京交通大学有关保留、使用学位论文的规定。特授权北京交通大学可以将学位论文的全部或部分内容编入有关数据库进行检索,提供阅览服务,并采用影印、缩印或扫描等复制手段保存、汇编以供查阅和借阅。同意学校向国家有关部门或机构送交论文的复印件和磁盘。学校可以为存在馆际合作关系的兄弟高校用户提供文献传递服务和交换服务。\par
  （保密的学位论文在解密后适用本授权说明）
}

% 学位论文数据集
\def\bjtu@label@dataset{学位论文数据集}

%%
%% labels in the english title page
%%
\def\bjtu@label@englishadvisor{Advisor}
\def\bjtu@label@englishcoadvisor{Co-advisor}
\def\bjtu@label@englishstatement{Submitted in total fulfillment
  of the requirements for the degree of \bjtu@value@englishdegree \\
  in \bjtu@value@englishmajor}

%%
%% labels in the abstracts
%%
\def\bjtu@label@introduction{序~~~~言}
\def\bjtu@label@abstract{摘~~~~要}
\def\bjtu@label@englishabstract{ABSTRACT}
\def\bjtu@label@keywords{关键词：}
\def\bjtu@label@englishkeywords{KEY WORDS:~}

%% labels in the summary
\def\bjtu@label@summary{结论}

%%
%% labels in the publications
%%
\def\bjtu@label@publications{攻读学位期间发表的学术论文}

%%
%% labels in the publications
%%
\def\bjtu@label@patents{攻读学位期间申请的专利}

%%
%% labels in the projects
%%
\def\bjtu@label@projects{攻读学位期间参与的项目}

%%
%% labels in the projects
%%
\ifbjtu@bachelor
  \def\bjtu@label@archivements{作者简历及攻读学士学位期间取得的研究成果}
\else
  \ifbjtu@master
    \def\bjtu@label@archivements{作者简历及攻读硕士学位期间取得的研究成果}
  \else
    \ifbjtu@doctor
      \def\bjtu@label@archivements{作者简历及攻读博士学位期间取得的研究成果}
    \else
      \ClassError{bjtuthesis}%
      {Unknown value for degree.}{}
    \fi 
  \fi
\fi


%%
%% labels in the resume
%%
\def\bjtu@label@resume{简~~~~历}


%%
%% labels in the thanks
%%
\def\bjtu@label@thanks{致~~~~谢}

%%%
%% listings name
%%
\def\bjtu@value@listingname{代码}

%%
%% the theorem name definitions
%%
\def\bjtu@label@algo{算法}
\def\bjtu@label@thm{定理}
\def\bjtu@label@lem{引理}
\def\bjtu@label@prop{命题}
\def\bjtu@label@cor{推论}
\def\bjtu@label@defn{定义}
\def\bjtu@label@conj{猜想}
\def\bjtu@label@exmp{例}
\def\bjtu@label@rem{注}
\def\bjtu@label@case{情形}
\def\bjtu@label@proof{证明}

\endinput

%% End of file `bjtuthesis.cfg'.

}


%==========
% Segment 1. Import LaTeX packages.
%==========

% tmargin = 16mm + 5mm = 21mm, bottom= 30mm + 1em(4mm), footskip = bottom - 15mm - 1em(4mm)
\RequirePackage[centering,a4paper,top=30mm,bottom=34mm,left=35mm,right=25mm,tmargin=21mm,footskip=11mm]{geometry} %设置版面

\RequirePackage{fancyhdr}
\RequirePackage{lastpage}
\RequirePackage{fontspec}
\RequirePackage{metalogo,doc}
\RequirePackage{threeparttable}
\RequirePackage{dcolumn}
\RequirePackage{multirow}
\RequirePackage{booktabs}
\RequirePackage{mathtools,amsthm,amsfonts,amssymb,bm,mathrsfs} 
\RequirePackage{upgreek}
\RequirePackage{graphicx}
\RequirePackage{subfigure}
\RequirePackage{longtable}
\RequirePackage{ccaption}
\RequirePackage[backend=biber, style=caspervector,utf8, bibencoding=utf8, sorting=none]{biblatex}
\RequirePackage{xcolor}
\RequirePackage{wasysym}
\RequirePackage{listings}
\RequirePackage[xetex, bookmarksnumbered, colorlinks, urlcolor=black, linkcolor=black, citecolor=black, plainpages=false, pdfstartview=FitH]{hyperref}
\RequirePackage{longtable}
\RequirePackage[perpage, bottom]{footmisc}
\RequirePackage[inline]{enumitem}
\RequirePackage{pdfpages}
\RequirePackage{calc}
\RequirePackage{algorithm, algorithmic}
\RequirePackage{siunitx}
\RequirePackage{tikz}
\usetikzlibrary{shapes.geometric, arrows}

%==========
% Segment 2. Define general-purpose LaTeX commands.
%==========

\setcounter{secnumdepth}{4}  % 章节编号深度 (part 对应 -1)
\setcounter{tocdepth}{2}     % 目录深度 (part 对应 -1)
    
% User defined command list
%% \me \mi \mj \dif \cleardoublepage \cndash \CJKLaTeX

% A new column type
\newcolumntype{d}[1]{D{.}{.}{#1}}% or D{.}{,}{#1} or D{.}{\cdot}{#1}

% upper math letter
\newcommand{\me}{\mathrm{e}} 
\newcommand{\mi}{\mathrm{i}}
\newcommand{\mj}{\mathrm{j}}
\newcommand{\dif}{\mathrm{d}}

\def\cleardoublepage{\clearpage%
	\if@openright%
	  \ifodd
	    \c@page%
	  \else%
  	    \thispagestyle{empty}%
  	    \hbox{}%
	    \newpage%
	  \fi%
	\fi}

% CJK-LaTeX Logo \CJKLaTeX
\newcommand{\CJKLaTeX}{CJK--\LaTeX} 

% cndash
\newcommand{\cndash}{\rule{0.0em}{0pt}\rule[0.35em]{1.4em}{0.05em}\rule{0.2em}{0pt}}

%==========
% Segment 3. Configure the imported packages, also extend LaTeX command in bjtuthesis
%==========

% Set the mainfont
\setmainfont{TeX Gyre Termes}
\setCJKmainfont[BoldFont=SimHei,ItalicFont=SimSun]{SimSun} % 强行指定斜体和粗体

%% 行距缩放因子
\linespread{1.3}

%% CTeX Settings
% 名称：目录、表格索引、插图索引、算法、算法索引、代码
\ctexset{contentsname={\bjtu@contentsname}}
\ctexset{listfigurename={\bjtu@listfigurename}}
%\ctexset{listtablename={\bjtu@listtablename}}
\floatname{algorithm}{\bjtu@label@algo}
\renewcommand{\listalgorithmname}{\bjtu@listalgorithmname}
\renewcommand{\lstlistingname}{\bjtu@value@listingname}
% Title Settings at the chapter Level
\ctexset{chapter={
  name={,},
  number={\arabic{chapter}},
  nameformat={\zihao{3}\bfseries},
  titleformat={\zihao{3}\bfseries},
	beforeskip={15\p@},
	afterskip={12\p@},
	}
}

\newcommand*\listoffiguresandtables{%
  \chapter*{\listfigureandtablename}% 如果是 article 文档类，把 \chapter 换成 \section
  \@mkboth{\MakeUppercase\listfigureandtablename}%
          {\MakeUppercase\listfigureandtablename}%

  \@starttoc{lof}%% 插图目录
  %\bigskip       %% 目录间的距离，不需要的话可以删掉或者换其他内容
  \@starttoc{lot}%% 表格目录
}

\newcommand*\listfigureandtablename{索引}

% Title Settings at the section Level
\ctexset{section={
	format={\zihao{-3}\bfseries},
	afterskip={1.0ex \@plus .2ex},
	}
}
% Title Settings at the subection Level
\ctexset{subsection={
  format={\zihao{4}\bfseries},
	indent={0\ccwd},
	afterskip={1.0ex \@plus .2ex},
	}
}
% Title Settings at the subsubection Level
\ctexset{subsubsection={
	format={\normalfont\normalsize},
	indent={0\ccwd},
	afterskip={1.0ex \@plus .2ex},
	}
}

% bullets in the item
\renewcommand{\labelitemi}{\ensuremath{\bullet}}

%% Graphic path & file extension
\graphicspath{{fig/}{figure/}{figures/}{logo/}{logos/}{graph/}{graphs}}
\DeclareGraphicsExtensions{.pdf,.eps,.png,.jpg,.jpeg}

% Caption setting in ccaption pacakge
\captiondelim{\ }
\captionnamefont{\songti\small}
\captiontitlefont{\songti\small}

% Floating parameters
\renewcommand{\textfraction}{0.15}
\renewcommand{\topfraction}{0.85}
\renewcommand{\bottomfraction}{0.65}
\renewcommand{\floatpagefraction}{0.60}

% 定公式、图、表、算法的编号为"3-1"的形式，即分隔符由.变为短杠
\renewcommand{\theequation}{\arabic{chapter}--\arabic{equation}}
\renewcommand{\thefigure}{\arabic{chapter}--\arabic{figure}}
\renewcommand\p@subfigure{\thefigure}
\renewcommand{\thetable}{\arabic{chapter}--\arabic{table}}
\renewcommand{\thealgorithm}{\arabic{chapter}--\arabic{algorithm}}

% 列表右缩进
\setlist[description]{leftmargin=\parindent,labelindent=\parindent}

% 流程图定义基本形状
\tikzstyle{startstop} = [rectangle, rounded corners, minimum width=2cm, minimum height=1cm, text centered, draw=black]
\tikzstyle{io} = [trapezium, trapezium left angle=75, trapezium right angle=105, minimum width=1cm, minimum height=1cm, text centered, draw=black]
\tikzstyle{process} = [rectangle, minimum width=2cm, minimum height=1cm, text centered, draw=black]
\tikzstyle{decision} = [diamond, minimum width=2cm, minimum height=1cm, text centered, draw=black]
\tikzstyle{arrow} = [thick,->,>=stealth]

% Setting Package listings/lstlistings
\lstset{tabsize=4, %
  frame=shadowbox, %把代码用带有阴影的框圈起来
  commentstyle=\color{red!50!green!50!blue!50},%浅灰色的注释
  rulesepcolor=\color{red!20!green!20!blue!20},%代码块边框为淡青色
  keywordstyle=\color{blue!90}\bfseries, %代码关键字的颜色为蓝色，粗体
  showstringspaces=false,%不显示代码字符串中间的空格标记
  stringstyle=\ttfamily, % 代码字符串的特殊格式
  keepspaces=true, %
  breakindent=22pt, %
  numbers=left,%左侧显示行号
  stepnumber=1,%
  numberstyle=\tiny, %行号字体用小号
  basicstyle={\footnotesize\ttfamily}, %
  showspaces=false, %
  flexiblecolumns=true, %
  breaklines=true, %对过长的代码自动换行
  breakautoindent=true,%
  breakindent=4em, %
  aboveskip=1em, %代码块边框
  %% added by http://bbs.ctex.org/viewthread.php?tid=53451
  fontadjust,
  captionpos=t,
  framextopmargin=2pt,framexbottommargin=2pt,abovecaptionskip=-3pt,belowcaptionskip=3pt,
  xleftmargin=4em,xrightmargin=4em, % 设定listing左右的空白
  texcl=true,
  % 设定中文冲突，断行，列模式，数学环境输入，listing数字的样式
  extendedchars=false,columns=flexible,mathescape=true
  numbersep=-1em
}

% 定理环境
\newtheoremstyle{break}% name
  {}%      Space above, empty = `usual value'
  {}%      Space below
  {\itshape}% Body font
  {}%         Indent amount (empty = no indent, \parindent = para indent)
  {\bfseries}% Thm head font
  {.}%        Punctuation after thm head
  {\newline}% Space after thm head: \newline = linebreak
  {}%         Thm head spec
\theoremstyle{plain}
  \newtheorem{algo}{\bjtu@label@algo~}[chapter]
  \newtheorem{thm}{\bjtu@label@thm~}[chapter]
  \newtheorem{lem}[thm]{\bjtu@label@lem~}
  \newtheorem{prop}[thm]{\bjtu@label@prop~}
  \newtheorem{cor}[thm]{\bjtu@label@cor~}
\theoremstyle{definition}
  \newtheorem{defn}{\bjtu@label@defn~}[chapter]
  \newtheorem{conj}{\bjtu@label@conj~}[chapter]
  \newtheorem{exmp}{\bjtu@label@exmp~}[chapter]
  \newtheorem{rem}{\bjtu@label@rem~}
  \newtheorem{case}{\bjtu@label@case~}
\theoremstyle{break}
  \newtheorem{bthm}[thm]{\bjtu@label@thm~}
  \newtheorem{blem}[thm]{\bjtu@label@lem~}
  \newtheorem{bprop}[thm]{\bjtu@label@prop~}
  \newtheorem{bcor}[thm]{\bjtu@label@cor~}
  \renewcommand{\proofname}{\bf\bjtu@label@proof}

%%%%%%%%%%%%%%%%%%%%%%%%%%%%%%%%%%%%%%%%%%%%%%%%%%%%%%%%%%%%%%%%
% The following definitions are to extend the LaTeX algorithmic 
% package with SWITCH statements and one-line structures.
% The extension is by 
%   Prof. Farn Wang 
%   Dept. of Electrical Engineering, 
%   National Taiwan University. 
% 
\newcommand{\SWITCH}[1]{\STATE \textbf{switch} (#1)}
\newcommand{\ENDSWITCH}{\STATE \textbf{end switch}}
\newcommand{\CASE}[1]{\STATE \textbf{case} #1\textbf{:} \begin{ALC@g}}
\newcommand{\ENDCASE}{\end{ALC@g}}
\newcommand{\CASELINE}[1]{\STATE \textbf{case} #1\textbf{:} }
\newcommand{\DEFAULT}{\STATE \textbf{default:} \begin{ALC@g}}
\newcommand{\ENDDEFAULT}{\end{ALC@g}}
\newcommand{\DEFAULTLINE}[1]{\STATE \textbf{default:} }
% 
% End of the LaTeX algorithmic package extension.
%%%%%%%%%%%%%%%%%%%%%%%%%%%%%%%%%%%%%%%%%%%%%%%%%%%%%%%%%%%%%%%%

%==========
% Segment 4. Draw the bjtuthesis
%==========

\newcommand\classification[1]{\def\bjtu@value@classification{#1}}
\newcommand\studentnumber[1]{\def\bjtu@value@studentnumber{#1}}
\newcommand\confidential[1]{\def\bjtu@value@confidential{#1}}
\newcommand\school[1]{\def\bjtu@value@school{#1}}
\newcommand\chinesedegree[1]{\def\bjtu@value@chinesedegree{#1}}
\renewcommand\title[2][\bjtu@value@title]{%
  \def\bjtu@value@title{#2}
  \def\bjtu@value@titlemark{\MakeUppercase{#1}}}
\renewcommand\author[1]{\def\bjtu@value@author{#1}}
\newcommand\advisor[1]{\def\bjtu@value@advisor{#1}}
\newcommand\advisorjobtitle[1]{\def\bjtu@value@advisorjobtitle{#1}}
\newcommand\coadvisor[1]{\def\bjtu@value@coadvisor{#1}}
\newcommand\major[1]{\def\bjtu@value@major{#1}}
\newcommand\submitdate[1]{\def\bjtu@value@submitdate{#1}}
\newcommand\defenddate[1]{\def\bjtu@value@defenddate{#1}}
\newcommand\institute[1]{\def\bjtu@value@institute{#1}}
\newcommand\chairman[1]{\def\bjtu@value@chairman{#1}}

% “绘制”BJTU中文封面
\renewcommand\makecover[1]{%
  \cleardoublepage
  \thispagestyle{empty}
  \begin{center}
  \vspace*{65.5pt}
  \includegraphics[width=0.72025\textwidth]{figure/bjtubanner}
  \vskip 30pt
  {\songti\zihao{2}\textbf{\renewcommand{\CJKglue}{\hskip 5pt plus 0.08\baselineskip}\bjtu@label@statement}}
  \vskip 60pt
  {\songti\zihao{-3}\bjtu@value@title}
  \vskip 15pt
  {\songti\zihao{-3}\bjtu@value@englishtitle}
  \vskip\stretch{1}
  {\songti\zihao{4}作者：\bjtu@value@author}
  \vskip 20pt
  {\songti\zihao{4}导师：\bjtu@value@advisor}
  \vskip 100pt
  {\songti\zihao{4}\bjtu@value@school}
  \vskip 20pt
  {\songti\zihao{4}\bjtu@value@defenddate}
  \end{center}
  \cleardoublepage
}


% “绘制”BJTU中文标题页
\renewcommand\maketitle[1]{%
  \cleardoublepage
  \thispagestyle{empty}
  \begin{center}
    \vspace*{4pt}
    \begin{tabular*}{\textwidth}{ @{}p{0.5\textwidth}  p{0.5\textwidth} }
      \zihao{5}学校代码：100004 & \hfill\zihao{5}密级：公开
    \end{tabular*}
  \end{center}
  \begin{center}
  \vspace*{25pt}
  {\kaishu\zihao{0}北京交通大学}
  \vskip 11pt
  {\songti\zihao{2}\textbf{\renewcommand{\CJKglue}{\hskip 5pt plus 0.08\baselineskip}\bjtu@label@statement}}
  \vskip 45pt
  {\songti\zihao{-3}\bjtu@value@title}
  \vskip 15pt
  {\songti\zihao{-3}\bjtu@value@englishtitle}
  \vskip 62pt
  \begin{tabular}{ @{}p{0.5\textwidth} p{0.5\textwidth}}
    \zihao{4}作者姓名：\bjtu@value@author & \zihao{4}学~~~~~~~~号：\bjtu@value@studentnumber \\
    & \\
    \zihao{4}导师姓名：\bjtu@value@advisor & \zihao{4}职~~~~~~~~称：\bjtu@value@advisorjobtitle \\
    & \\
    \zihao{4}
      \ifbjtu@professional\relax
        工程\bjtu@value@chinesedegree 专业领域：\bjtu@value@major
      \else
        \bjtu@value@chinesedegree 专业领域：\bjtu@value@major
      \fi
       & \zihao{4}学位级别：\bjtu@value@chinesedegree  \\
    & \\
  \end{tabular}
  \vskip\stretch{1}
  {\fangsong\zihao{4}}
  \def\tabcolsep{1pt}
  \def\arraystretch{1.5}
  {\songti\zihao{4}\bjtu@value@school}
  \vskip 20pt
  {\songti\zihao{4}\bjtu@value@defenddate}
  \end{center}
  \cleardoublepage
}


% English Title Page
\newcommand\englishtitle[1]{\def\bjtu@value@englishtitle{#1}}
\newcommand\englishauthor[1]{\def\bjtu@value@englishauthor{#1}}
\newcommand\englishadvisor[1]{\def\bjtu@value@englishadvisor{#1}}
\newcommand\englishcoadvisor[1]{\def\bjtu@value@englishcoadvisor{#1}}
\newcommand\englishschool[1]{\def\bjtu@value@englishschool{#1}}
\newcommand\englishinstitute[1]{\def\bjtu@value@englishinstitute{#1}}
\newcommand\englishdate[1]{\def\bjtu@value@englishdate{#1}}
\newcommand\englishdegree[1]{\def\bjtu@value@englishdegree{#1}}
\newcommand\englishmajor[1]{\def\bjtu@value@englishmajor{#1}}
% "绘制"英文标题页
\newcommand\makeenglishtitle[1]{%
  \cleardoublepage
  \thispagestyle{empty}
  \begin{center}
      {\normalfont\zihao{4} \bjtu@label@englishstatement} 
      \vspace{20pt} \vskip\stretch{1}
      {\huge\bjtu@value@englishtitle \vskip 1pt}
      \vskip \stretch{1}
    \ifbjtu@review
      \vskip \stretch{1}
      \vskip 3pt
      \vskip \stretch{2}
    \else
      {\normalfont\zihao{4}\bjtu@value@englishauthor} 
      \vskip \stretch{1}
      {\normalfont\zihao{4}\bjtu@label@englishadvisor} 
      \vskip 2pt
      {\normalfont\zihao{4}\bjtu@value@englishadvisor} 
      \vskip 5pt
      \ifx\bjtu@value@englishcoadvisor\undefined\else
        {\normalfont\zihao{4}\bjtu@label@englishcoadvisor} 
        \vskip 2pt
        {\normalfont\zihao{4}\bjtu@value@englishcoadvisor} 
        \vskip \stretch{2}
      \fi
    \fi 
    \normalfont\bjtu@value@englishinstitute \vskip 30pt
    \normalfont\bjtu@value@englishdate \vskip 20pt
  \end{center}
  \cleardoublepage
}

% 独创性声明
\newcommand\makeDeclareOriginal{%
  \cleardoublepage
  \chapter*{\bjtu@label@original\markboth{\bjtu@label@original}{}}
  {\par\zihao{-4}\bjtu@label@originalcontent\par}
  \vskip 40pt
  \hspace{2em}
  \bjtu@label@authorsign
  \hfill
  \hfill
  \hfill
  \bjtu@label@originalDate
  \hspace{1.5em}
  \cleardoublepage
}


% 授权声明
\newcommand\makeDeclareAuthorization{%
  \cleardoublepage
  \thispagestyle{empty}
  \begin{center}
    \vspace*{35pt}
  	{\bf\zihao{3} \bjtu@label@authorization}
  \end{center}
  \vskip 10pt
  {\par\zihao{-4}\bjtu@label@authorizationcontent\par}
  \vskip 105pt
  \par\bjtu@label@authorsign\hfill\bjtu@label@Supervisorsign\hfill\hspace{2em}\par
  \vskip 15pt
  \par\bjtu@label@originalDate\hfill\hspace{2em}\bjtu@label@originalDate\par
  \cleardoublepage
}

% 学位论文数据集
\newcommand\makeDataset{%
  \cleardoublepage
  \chapter*{\bjtu@label@dataset\markboth{\bjtu@label@dataset}{}}
  %\vskip 10pt
  \begin{center}
    \begin{table}[!hpb]
      \centering
      \caption{数据集页}
      \begin{tabular}{ | p{7em} | p{4em} | p{9em} | p{8em} | p{6em} |}
        \hline
        关键词* & 密级* & 中图分类号 & UDC & 论文资助 \\ \hline
          &  &  &  &  \\ \hline
        \multicolumn{2}{|l|}{学位授予单位名称*} & 学位授予单位代码* & 学位类别* & 学位级别* \\ \hline
        \multicolumn{2}{|l|}{北京交通大学} & 10004 &  &  \\ \hline
        \multicolumn{2}{|l|}{论文题名*} & \multicolumn{2}{l|}{并列题名} & 论文语种* \\ \hline
        \multicolumn{2}{|l|}{} & \multicolumn{2}{l|}{} &  \\ \hline
        作者姓名* & \multicolumn{2}{l|}{} & 学号* &  \\ \hline
        \multicolumn{2}{|l|}{培养单位名称*} & 培养单位代码* & 培养单位地址 & 邮编 \\ \hline
        \multicolumn{2}{|l|}{北京交通大学} & 10004 & 北京市海淀区西直门外上园村3号 & 100044 \\ \hline
        \multicolumn{2}{|l|}{工程领域*} & 研究方向* & 学制* & 学位授予年* \\ \hline
        \multicolumn{2}{|l|}{} &  &  &  \\ \hline
        论文提交日期* & \multicolumn{4}{l|}{} \\ \hline
        论导师姓名* & \multicolumn{2}{l|}{} & 职称* & \\ \hline
        评阅人 & \multicolumn{2}{l|}{答辩委员会主席*} & \multicolumn{2}{l|}{答辩委员会成员} \\ \hline
         & \multicolumn{2}{l|}{} & \multicolumn{2}{l|}{} \\
         & \multicolumn{2}{l|}{} & \multicolumn{2}{l|}{} \\ \hline
        \multicolumn{5}{|l|}{电子版论文提交格式  文本（ ）  图像（ ） 视频（ ） 音频（ ） 多媒体（ ） 其他（ ）} \\
        \multicolumn{5}{|l|}{推荐格式：application/msword；application/pdf} \\ \hline
        \multicolumn{2}{|l|}{电子版论文出版（发布）者} & \multicolumn{2}{l|}{电子版论文出版（发布）地} & 权限声明 \\ \hline
        \multicolumn{2}{|l|}{} & \multicolumn{2}{l|}{} &  \\ \hline
        论文总页数* & \multicolumn{4}{l|}{} \\ \hline
        \multicolumn{5}{|l|}{共33项，其中带*为必填数据，为21项。} \\ \hline
      \end{tabular}
    \end{table}
  \end{center}
  %\cleardoublepage
}

% fancyhdr页眉页脚设置

% 不显示页眉中的编号
\pagestyle{fancy}\renewcommand{\chaptermark}[1]{\markboth{\MakeUppercase{#1}}{}}

\ifbjtu@bachelor
	% 本科学位论文页眉页脚设置
	%% 正文页眉页脚
	\fancypagestyle{main}{
	  \fancyhf{}
	  \fancyhead[L]{\includegraphics[width=0.22\textwidth]{figure/bjtubanner}}
	  \fancyhead[R]{\nouppercase{\small\songti\bjtu@value@titlemark}}
	  \fancyfoot[C]{\small 第~{\bfseries\thepage}~页\,共~{\bfseries\pageref{LastPage}}~页}
	  \renewcommand{\headheight}{32pt}
	}
	%% 开章页页眉页脚
	\fancypagestyle{plain}{% 设置开章页页眉页脚风格(只有页码作为页脚)
	  \fancyhf{}%
	  \fancyhead[L]{\nouppercase{\small\songti\includegraphics[width=0.22\textwidth]{figure/bjtubanner}}}
	  \fancyhead[R]{\nouppercase{\small\songti\bjtu@value@titlemark}}
	  \fancyfoot[C]{\small 第~{\bf\thepage}~页\,共~{\bf\pageref{LastPage}}~页}
	  \renewcommand{\headheight}{32pt}
	}
\else
	% 研究生学位论文页眉页脚设置
	\if@twoside
	  %% 正文页眉页脚
	  \fancypagestyle{main}{
	    \fancyhead[CE]{
        \small\songti\bjtu@value@school
        \bjtu@value@chinesedegree
        % 专硕
        \ifbjtu@professional\relax
          专业
        \fi
        \bjtu@label@thesis
      }
      \fancyhead[L]{}
      \fancyhead[R]{}
      \fancyhead[CO]{\nouppercase{\small\songti\leftmark}}
	    \fancyfoot[C]{\small ~{\bfseries\thepage}~}
	    \renewcommand{\headheight}{32pt}
	  }
	  %% 开章页页眉页脚
	  \fancypagestyle{plain}{
	    \fancyhf{}
	    \fancyhead[CE]{\nouppercase{
        \small\songti\bjtu@value@school
        \bjtu@value@chinesedegree
        % 专硕
        \ifbjtu@professional\relax
          专业
        \fi
        \bjtu@label@thesis
      }}
      \fancyhead[L]{}
      \fancyhead[R]{}
	    \fancyhead[CO]{\nouppercase{\small\songti\leftmark}}
	    \fancyfoot[C]{\small ~{\bf\thepage}~}
	    \renewcommand{\headheight}{32pt}
	  }
	\else
	%% 正文页
	\fancypagestyle{main}{
	  \fancyhead[L]{
      \small\songti\bjtu@value@school
      \bjtu@value@chinesedegree
      % 专硕
      \ifbjtu@professional\relax
        专业
      \fi
      \bjtu@label@thesis
    }
	  \fancyhead[R]{\nouppercase{\small\songti\leftmark}}
	  \fancyfoot[C]{\small ~{\bf\thepage}~}
	  \renewcommand{\headheight}{32pt}
	}
	\fancypagestyle{plain}{
	  \fancyhf{}
	  \fancyhead[L]{\nouppercase{
      \small\songti\bjtu@value@school
      \bjtu@value@chinesedegree
      % 专硕
      \ifbjtu@professional\relax
        专业
      \fi
      \bjtu@label@thesis
    }}
    \fancyhead[R]{\nouppercase{\small\songti\leftmark}}
	  % \fancyhead[L]{\small {\it\bjtu@value@titlemark}}
	  \fancyfoot[C]{\small ~{\bf\thepage}~}
	  \renewcommand{\headheight}{32pt}
	}
	\fi
\fi

% 序言
\newenvironment{introduction}{
  \chapter*{\bjtu@label@introduction}
  \markboth{\bjtu@label@introduction}
  {}
}

% 中英文摘要
\newenvironment{abstract}{\chapter*{\bjtu@label@abstract}\markboth{\bjtu@label@abstract}{}}{}
\newcommand\keywords[1]{\noindent{\bf\songti\zihao{-4} \bjtu@label@keywords} #1}
\newenvironment{englishabstract}{\chapter*{\bfseries \bjtu@label@englishabstract}\markboth{\bjtu@label@englishabstract}{}}{}
\newcommand\englishkeywords[1]{\noindent{\bf\zihao{-4} \bjtu@label@englishkeywords} #1}

\renewcommand\tableofcontents{%
  \chapter*{\contentsname}
  \@mkboth{\MakeUppercase\contentsname}{\MakeUppercase\contentsname}%
  \@starttoc{toc}%
}

\addtocontents{toc}{\let\string\CTEX@spaceChar\relax}

\newenvironment{publications}[1]
     {\chapter{\bjtu@label@publications}%
      \@mkboth{\MakeUppercase\bjtu@label@publications}
              {\MakeUppercase\bjtu@label@publications}%
      \list{\@biblabel{\@arabic\c@enumiv}}%
           {\settowidth\labelwidth{\@biblabel{#1}}%
            \leftmargin\labelwidth
            \advance\leftmargin\labelsep
            \@openbib@code
            \usecounter{enumiv}%
            \let\p@enumiv\@empty
            \renewcommand\theenumiv{\@arabic\c@enumiv}}%
      \sloppy
      \clubpenalty4000
      \@clubpenalty \clubpenalty
      \widowpenalty4000%
      \sfcode`\.\@m}
     {\def\@noitemerr
       {\@latex@warning{Empty `publications' environment}}%
      \endlist}

\newenvironment{patents}[1]
     {\chapter{\bjtu@label@patents}%
      \@mkboth{\MakeUppercase\bjtu@label@patents}
              {\MakeUppercase\bjtu@label@patents}%
      \list{\@biblabel{\@arabic\c@enumiv}}%
           {\settowidth\labelwidth{\@biblabel{#1}}%
            \leftmargin\labelwidth
            \advance\leftmargin\labelsep
            \@openbib@code
            \usecounter{enumiv}%
            \let\p@enumiv\@empty
            \renewcommand\theenumiv{\@arabic\c@enumiv}}%
      \sloppy
      \clubpenalty4000
      \@clubpenalty \clubpenalty
      \widowpenalty4000%
      \sfcode`\.\@m}
     {\def\@noitemerr
       {\@latex@warning{Empty `patents' environment}}%
      \endlist}

\newenvironment{archivements}[1]
     {\chapter{\bjtu@label@archivements}%
      \@mkboth{\MakeUppercase\bjtu@label@archivements}
              {\MakeUppercase\bjtu@label@archivements}%
      \list{\@biblabel{\@arabic\c@enumiv}}%
           {\settowidth\labelwidth{\@biblabel{#1}}%
            \leftmargin\labelwidth
            \advance\leftmargin\labelsep
            \@openbib@code
            \usecounter{enumiv}%
            \let\p@enumiv\@empty
            \renewcommand\theenumiv{\@arabic\c@enumiv}}%
      \sloppy
      \clubpenalty4000
      \@clubpenalty \clubpenalty
      \widowpenalty4000%
      \sfcode`\.\@m}
     {\def\@noitemerr
       {\@latex@warning{Empty `archivements' environment}}%
      \endlist}

\newenvironment{projects}[1]
     {\chapter{\bjtu@label@projects}%
      \@mkboth{\MakeUppercase\bjtu@label@projects}
              {\MakeUppercase\bjtu@label@projects}%
      \list{\@biblabel{\@arabic\c@enumiv}}%
           {\settowidth\labelwidth{\@biblabel{#1}}%
            \leftmargin\labelwidth
            \advance\leftmargin\labelsep
            \@openbib@code
            \usecounter{enumiv}%
            \let\p@enumiv\@empty
            \renewcommand\theenumiv{\@arabic\c@enumiv}}%
      \sloppy
      \clubpenalty4000
      \@clubpenalty \clubpenalty
      \widowpenalty4000%
      \sfcode`\.\@m}
     {\def\@noitemerr
       {\@latex@warning{Empty `projects' environment}}%
      \endlist}

\newenvironment{resume}
  {\chapter{\bjtu@label@resume}}
  {}
  
\newenvironment{resumesection}[1]
  {{\noindent\normalfont\bfseries #1}
   \list{}{\labelwidth\z@
           \leftmargin 2\ccwd}
   \item\relax}
   {\endlist}
   
\newenvironment{resumelist}[1]
  {{\noindent\normalfont\bfseries #1}
   \list{}{\labelwidth\z@
           \leftmargin 4\ccwd
           \itemindent -2\ccwd
           \listparindent\itemindent}
   \item\relax}
   {\endlist}
   
\renewenvironment{thanks}{%
  \cleardoublepage
  \thispagestyle{empty}
  \begin{center}
    \vspace*{35pt}
    {\heiti\zihao{3}\bjtu@label@thanks}
  \end{center}
  {}
}

\newenvironment{summary}
{\chapter*{\bjtu@label@summary\markboth{\bjtu@label@summary}{}}\addcontentsline{toc}{chapter}{\bjtu@label@summary}
}{}

\newenvironment{symbolnote}{\chapter{\bjtu@label@symbolnote}}{} 

\endinput

%%
%% End of file `bjtuthesis.cls'.



